\section{Synthetic Counting}
\label{sec:synth-counting}

\begin{keyterm}[Counting Number]
\label{def:counting-number}
A \textbf{counting number} is any positive whole integer used to enumerate
objects: $1, 2, 3, \ldots$ and so on.
\end{keyterm}

\begin{property}[Successor Rule]
\label{prop:successor}
For any counting number $n$, the next counting number is $n + 1$. This
process never terminates.
\end{property}

\begin{example}[Listing Successors]
\label{ex:succ-list}
Starting at $n = 5$, the next three counting numbers are $6$, $7$, and $8$.
\end{example}

\begin{checkpoint}[Try Counting]
\label{cp:try-count}
List the next five counting numbers after $12$. Answer: $13, 14, 15, 16, 17$.
\end{checkpoint}

\begin{note}[Historical]
\label{note:origin}
Counting systems were among the earliest mathematical innovations; every
civilization eventually invents them in some form.
\end{note}
